\chapter{引言}

\section{概述}

随着现代科学技术的发展和人们生活质量的提高，对于工程结构的性能提出
了越来越高的要求。例如：现代精密仪器、大型设备往往对于振动与位移有严格
的限制；生命线工程结构，要求在大震和大灾作用下依然保有必要的功能，以为
灾后救援与重建提供保障。20 世纪中叶以来，尽管社会发展水平有了巨大的提
高，然而由于灾害性作用而造成的损失却反而越来越大，这给结构工程学科带来
了一系列新的挑战性课题。正是在这样的背景下，基于性能的设计思想开始浮出
水面，并在近十年来引起了学者们强烈的兴趣


``自然界只有一个，自然现象遵循着不依赖于人类意志的客观规律。然而，
数理科学中却有着两套反映这些规律的体系：确定性描述和概率论描述。''虽
然概率论方法的发展引起了科学家和哲学家们关于自然本质的讨论，但是直到本
世纪五十年代以前，两套方法在各自独立的领域内都得到了长足的发展。六十年
代以来，由于本质非线性行为特别是混沌、分形等现象的发现和深入研究，随机
方法的重要性得到了日益深刻的认识。人们发现，在确定性非线性系统的长期
演化行为中会出现与随机行为不能加以区别的现象。而采用概率密度演化描述的
方法却能很好地描述其演化密度的长期行为\cite{chen2021,soave2020}。

\section{研究现状}

这里综述相关领域的研究现状，并说明当前研究中仍待解决的问题。

\subsection{理论研究}

三级标题采用黑体四号。公式居中，编号右对齐：
\begin{equation}
  E = mc^2 .
\end{equation}

图应清晰可辨，图题置于图下方；表题通常置于表上方。

\begin{figure}[htbp]
  \centering
  \fbox{\rule{0pt}{3.5cm}\rule{8cm}{0pt}}
  \caption{示例图题}
  \label{fig:example}
\end{figure}

\begin{table}[htbp]
  \centering
  \caption{示例表题}
  \label{tab:example}
  \begin{tabular}{ccc}
    \toprule
    项目 & 数值一 & 数值二 \\
    \midrule
    示例 & 1.0 & 2.0 \\
    \bottomrule
  \end{tabular}
\end{table}
